\section{Preliminaries}
\label{sec:preliminaries}

We collect the approximation, comparison and diagonalization tools for
the dimension bound, including comparison from a uniform rank gap for
all lower semicontinuous extended traces. The central-sequence inputs
for the independent absorption implication are stated in
Section~\ref{sec:central}.

All $C^*$-algebras are complex. An unadorned tensor product is spatial;
maximal tensor products are explicitly marked $\otimes_{\max}$.
Representations are nondegenerate.  We write $D_+$ for the positive cone of
a $C^*$-algebra $D$, $\widetilde D$ for its unitization, and
\[
 \her_D(h)=\overline{hDh}\quad(h\in D_+).
\]
For any subset $S\subset D$, let $\Ideal_D(S)$ denote the closed two-sided
ideal generated by $S$; for $x\in D$ we abbreviate
$\Ideal_D(x)=\Ideal_D(\{x\})$.  The ambient algebra is omitted when it
is clear.  For $\varepsilon>0$,
$(h-\varepsilon)_+$ denotes the usual spectral cut.  A full element is
one generating the whole ambient algebra as a closed ideal.  A nonzero
$h\in D_+$ is \emph{totally full} if $f(h)$ is full for every nonzero
$f\in C_0((0,\|h\|])_+$.  For a subset $S$ of an ambient algebra,
$S^\perp$ consists of the elements annihilating $S$ on both sides.

\subsection{Order zero and finite-dimensional approximation}

A completely positive contractive map is abbreviated to c.p.c.  A
completely positive map $\sigma:D\to E$ has order zero if
$\sigma(a)\sigma(b)=0$ whenever $a,b\in D_+$ and $ab=0$.
For a c.p.c. order-zero map $\sigma$, the structure theorem of Winter--Zacharias
\cite[Theorem~3.3 and Corollaries~4.1--4.2]{WZOrderZero} gives a
$*$-homomorphism $\sigma_{\mathrm{hom}}$ and a commuting positive
contraction $h$ in the multiplier algebra of $C^*(\sigma(D))$ such that
\[
 \sigma(a)=h\sigma_{\mathrm{hom}}(a).
\]
For $f\in C_0((0,1])_+$, the formula
$f(\sigma)(a)=f(h)\sigma_{\mathrm{hom}}(a)$ defines an order-zero map
into $E$, contractive when $\|f\|\le1$.  We write
$\sigma^{1/l}$ for $f(t)=t^{1/l}$.  In particular,
\begin{equation}
 \label{prelim:order-zero}
 \|\sigma(a)\sigma(b)\|\le\|ab\|\quad(a,b\in D),
 \qquad
 \sigma(a)\sigma(b)=\sigma(1)\sigma(ab)
 \quad\text{if $D$ is unital}.
\end{equation}
The first inequality follows by writing the product as
$h^2\sigma_{\mathrm{hom}}(ab)$.  Multiplication of a c.p.c. order-zero map by
a positive contraction commuting with its entire range again gives a
c.p.c. order-zero map.

\Needspace{13\baselineskip}
\begin{definition}[Winter--Zacharias {\cite[Definition~2.1]{WZNuclear}}]
\label{prelim:nuclear-dimension}
For an integer $d\ge0$, a $C^*$-algebra $D$ has nuclear dimension at
most $d$, written $\nucdim(D)\le d$, if for every finite
$\mathcal F\subset D$ and $\varepsilon>0$ there are finite-dimensional
algebras $E_0,\ldots,E_d$, a c.p.c. map
$\psi:D\to\bigoplus_{j=0}^dE_j$, and c.p.c. order-zero maps
$\varphi_j:E_j\to D$ such that
\[
 \left\|\sum_{j=0}^d\varphi_j\psi_j(a)-a\right\|<\varepsilon
 \qquad(a\in\mathcal F).
\]
The $d+1$ summands are called colours.  The sum of the return maps
need not be contractive; each coloured return map must be contractive.
\end{definition}

We use the following approximation theorem in hereditary subalgebras as
well as in unital algebras.

\begin{theorem}[Brown--Carri\'on--White]
\label{prelim:bcw}
For every separable nuclear $D$ there are finite-dimensional algebras
$F_n$ and c.p.c. maps
\[
 D\xrightarrow{\alpha_n}F_n\xrightarrow{\beta_n}D
\]
such that $\beta_n\alpha_n\to\id_D$ in point-norm,
$\|\alpha_n(a)\alpha_n(b)\|\to0$ for orthogonal $a,b\in D_+$, and
each $\beta_n$ is a finite convex combination of c.p.c. order-zero maps
$F_n\to D$.
\end{theorem}

The convex order-zero decomposition is due to
Hirshberg--Kirchberg--White \cite[Theorem~1.4]{HKW}; the asymptotic
orthogonality refinement used here is \cite[Theorem~3.1]{BCW}.
In particular, the map induced by the
$\alpha_n$ in $\prod_nF_n/\bigoplus_nF_n$ is order zero, so
\begin{equation}
 \label{prelim:asymptotic-product}
 \limsup_n\|\alpha_n(a)\alpha_n(b)\|\le\|ab\|
 \qquad(a,b\in D).
\end{equation}
Nuclearity passes to hereditary subalgebras and
to unitizations.  A c.p.c. map from a separable nuclear algebra to a
quotient has a c.p.c. lift by the Choi--Effros lifting theorem
\cite[Theorem~3.10]{ChoiEffros}.  Every completely positive map from a nuclear algebra is
nuclear, by composition with finite-dimensional approximations of its
domain.

\subsection{Comparison and extended traces}

Write $\mathcal K$ for the compact operators on a separable infinite-dimensional
Hilbert space. For $u,v\in(D\otimes\mathcal K)_+$, write $u\precsim v$ if
$s_n^*vs_n\to u$ in norm for some $s_n\in D\otimes\mathcal K$.
Mutual comparison is Cuntz equivalence.  The semigroup $\Cu(D)$ consists
of the resulting classes, with block-sum addition and the order induced
by $\precsim$.  Its subsemigroup $W(D)$ consists of classes represented
in finite matrices over $D$.  We use the following standard properties:
\[
 [u+v]\le[u]+[v],\qquad
 [u+v]=[u]+[v]\ \text{if }uv=0,\qquad
 [s^*s]=[ss^*].
\]
Positive domination implies Cuntz comparison; elements of $\her(v)_+$
are Cuntz below $v$.  If $u\in\Ideal(v)_+$, then, for every
$\varepsilon>0$, there is $N<\infty$ with
$[(u-\varepsilon)_+]\le N[v]$.  Also
\[
 [u]=\sup_{\varepsilon>0}[(u-\varepsilon)_+],
 \qquad [(u-\varepsilon)_+]\ll[u],
\]
where $x\ll y$ means that any increasing sequence with supremum at
least $y$ has a term dominating $x$.  Spectral cuts are cofinal below
$[u]$ for this relation.  Projection classes are compact, and if a
projection $p$ is Cuntz below $v$, then $\her(v)$ contains a projection
Murray--von Neumann equivalent to $p$ (with matrix amplification when
needed).  These properties hold without simplicity or unitality; see
\cite[Section~4]{Rordam} and \cite[Sections~2 and~4]{ERS}.

\begin{lemma}[Finite compressions of a full element]
\label{prelim:full-compressions}
If $v\in D_+$ is full and $D$ is unital, then there are
$y_1,\ldots,y_N\in D$ such that
\[
 \sum_{l=1}^N y_l^*vy_l=1_D.
\]
\end{lemma}

\begin{proof}
The ideal criterion above, applied to the projection $1_D$, gives
$[1_D]\le N[v]$ for some $N$.  A sufficiently accurate Cuntz
approximation, compressed to the $1_D$ corner, gives
$x_1,\ldots,x_N\in D$ with
$\|1_D-\sum_lx_l^*vx_l\|<1$.  The positive element
$z=\sum_lx_l^*vx_l$ is invertible.  Taking $y_l=x_lz^{-1/2}$
gives the identity.
\end{proof}

A dimension function on $W(D)$ is an additive order-preserving map
$d:W(D)\to[0,\infty]$ with $d(0)=0$; no continuity is implicit.
The semigroup is \emph{almost unperforated} when
$(m+1)x\le my$ implies $x\le y$, for $m\in\N$.

\begin{theorem}[R\o rdam's comparison criterion]
\label{prelim:comparison}
If $D\cong D\otimes\Z$, then $W(D)$ is almost unperforated.
More generally, if $W(D)$ is almost unperforated, $u,v\in D_+$,
$u\in\Ideal_D(v)$, and
\[
 d([u])<d([v])\quad\text{for every dimension function $d$ with }d([v])=1,
\]
then $u\precsim v$.
\end{theorem}

These are \cite[Theorem~4.5 and Corollary~4.7]{Rordam}.  The second
assertion uses abstract dimension functions, and requires neither
exactness nor a trace-space hypothesis.

A lower semicontinuous extended trace on $D$ is an additive homogeneous
map $\tau:D_+\to[0,\infty]$ satisfying
$\tau(x^*x)=\tau(xx^*)$ and norm lower semicontinuity.
We allow traces that are not densely finite.  Their extensions to finite
matrices use the unnormalized matrix trace, and their extensions to the
stabilization are understood.  For a positive contraction $z$, put
\[
 d_\tau(z)=\lim_{l\to\infty}\tau(z^{1/l}).
\]
Scaling a nonzero positive element to a contraction gives the same
definition for arbitrary $z\ge0$.  We use
$\tau(z)\le d_\tau(z)$ for contractions and
\begin{equation}
 \label{prelim:trace-cut}
 d_\tau((z-\eta)_+)\le\eta^{-1}\tau(z)\qquad(\eta>0).
\end{equation}

The next lemma follows from Elliott--Robert--Santiago's comparison
criterion \cite[Proposition~6.2]{ERS} in the presence of a uniform rank
gap.  If $\Cu(D)$ is almost unperforated, their criterion gives
$x\le y$ for $x,y\in\Cu(D)$ from the ideal containment
$x\le\sup_n ny$ and the strict inequalities
$\lambda(x)<\lambda(y)$ for every functional $\lambda$ with
$\lambda(y)=1$.  A functional is an additive order-preserving map to
$[0,\infty]$ that vanishes at zero and preserves increasing suprema.
For exact algebras, the extended trace correspondence converts the rank
gap into these inequalities, while the trace that is zero on an ideal
and infinite outside it detects ideal containment.  We give the direct
regularization proof.

\begin{lemma}[Comparison with a uniform rank gap]
\label{prelim:half-comparison}
Let $D$ be exact and $\Z$-stable.  If $u,v\in D_+$ satisfy
\[
 d_\tau(u)\le\tfrac12d_\tau(v)
 \quad\text{for every lower semicontinuous extended trace $\tau$ on $D$},
\]
then $u\precsim v$.
\end{lemma}

\begin{proof}
Let $I=\Ideal_D(v)$.  The map that is zero on $I_+$ and infinity on
$D_+\setminus I$ is a lower semicontinuous extended trace: additivity
follows from heredity of $I$, the trace identity from
$x^*x\in I\Longleftrightarrow xx^*\in I$, and lower semicontinuity
from closedness of $I$.  The assumed inequality for this trace forces
$u\in I$.

Work in $D\otimes\mathcal K$, whose $W$-semigroup identifies with
$\Cu(D)$ and is almost unperforated.  Fix a dimension function $d$ on
this semigroup with $d([v])=1$.  Its regularization is
\[
 \overline d([z])=\sup_{\delta>0}d([(z-\delta)_+]).
\]
By \cite[Lemma~4.7]{ERS} and the spectral-cut description of compact
containment, $\overline d$ is a functional on $\Cu(D)$.  The
correspondence in \cite[Proposition~4.2]{ERS} realizes it as
$d_\lambda$ for a lower semicontinuous extended $2$-quasitrace.
Haagerup proved that finite-valued $2$-quasitraces on exact unital
algebras are traces.  Blanchard--Kirchberg extended this conclusion to
lower semicontinuous $2$-quasitraces on exact algebras, allowing a
possibly proper finite-domain ideal \cite[Remarks~4.3 and~4.5]{ERS}.
Thus $\lambda$ is a trace.  Restrict $\lambda$ to the original corner
$D$.  For each fixed $\eta>0$ the hypothesis gives
\[
 \begin{split}
 d([(u-\eta)_+])
 &\le\overline d([u])=d_\lambda(u)
 \le\tfrac12d_\lambda(v)
 =\tfrac12\overline d([v])\\
 &\le\tfrac12<1=d([v]).
 \end{split}
\]
The strict inequality is against $d([v])$, so it remains valid when
$\overline d([v])=0$.  Theorem~\ref{prelim:comparison} now compares
every $(u-\eta)_+$ with $v$.  Taking the supremum over cuts gives
$u\precsim v$.  Finally, comparisons between elements of the original
corner can be witnessed in that corner, by compressing each witness on
both sides by its multiplier projection.
\end{proof}

\begin{lemma}[Extended traces on finite matrices]
\label{prelim:matrix-traces}
Let $A$ be unital, let $k\ge1$, and let $\lambda$ be a lower
semicontinuous extended trace on $A\otimes M_k$.  Set
$\tau(a)=\lambda(a\otimes1_k)$ and
$L_k=\id_A\otimes\operatorname{tr}_k$, with normalized matrix trace.
Then
\[
 \lambda(z)=\tau(L_k(z))\qquad(z\in(A\otimes M_k)_+).
\]
No finiteness assumption on $\lambda$ is needed.
\end{lemma}

\begin{proof}
Write $e_{ij}=1_A\otimes e_{ij}$.  Additivity and the trace identity give
\[
 \lambda(z)
 =\sum_{i=1}^k\lambda(z^{1/2}e_{ii}z^{1/2})
 =\sum_{i=1}^k\lambda(e_{ii}ze_{ii}).
\]
Unitary equivalence of the diagonal corners gives
$\lambda(a\otimes e_{ii})=k^{-1}\tau(a)$ for $a\in A_+$.
Summing proves the formula.  Only finite sums and multiplication by
positive scalars occur, so the argument also applies to infinite values.
\end{proof}

\subsection{Ultraproducts and bounded witnesses}

Fix a free ultrafilter $\omega$ on $\N$.  For a sequence $(D_n)$ we use
the norm ultraproduct
\[
 \prod_\omega D_n
 =\prod_nD_n\big/\{(x_n):\lim_{n\to\omega}\|x_n\|=0\}.
\]
For a constant sequence this is denoted $D_\omega$; the ordinary
sequence algebra is $D_\infty=\ell^\infty(D)/c_0(D)$.

The following form of Kirchberg's $\varepsilon$-test is recorded in
\cite[Lemma~1.10]{BBSTWW}; we include the diagonal proof.

\begin{lemma}[Kirchberg's $\varepsilon$-test]
\label{prelim:epsilon-test}
Let $X_n$ be nonempty sets and let
$f_{j,n}:X_n\to[0,\infty)$ be given for $j,n\in\N$.
Suppose that for every $m$ and $\delta>0$ there is
$(x_n)\in\prod_nX_n$ with
$\lim_{n\to\omega}f_{j,n}(x_n)<\delta$ for $1\le j\le m$.
Then there is a single $(x_n)$ for which all these ultralimits are zero.
\end{lemma}

\begin{proof}
Choose $(x_n^{(m)})$ solving the first $m$ conditions with ultralimits
less than $1/(2m)$.  Choose $A_m\in\omega$ on which the first $m$
errors of this choice are less than $1/m$, and set
$E_m=A_1\cap\cdots\cap A_m$.  These sets are decreasing and belong
to $\omega$.
At coordinate $n$, use $x_n^{(l(n))}$, where
$l(n)=\max\{m\le n:n\in E_m\}$, using any element of $X_n$ if the
set is empty.  Then $l(n)\to_\omega\infty$, and every fixed error tends
to zero.
\end{proof}

We apply this to countably many norm relations on contractions and on
variables with fixed, separately prescribed norm bounds.  Nuclear
c.p.c. lifting permits the same procedure for maps out of a separable
nuclear algebra; a countable dense set records linear relations,
complete positivity and the needed order-zero relations.  Countably
many extra commutators or norms after a fixed contractive map can be
included among the tests.  The following elementary estimate supplies
the bounds for Cuntz witnesses.

\begin{lemma}[A spare spectral cut gives bounded witnesses]
\label{prelim:bounded-witness}
Let $u,y\in D_+$ and $0\le a<b$.  If
$u\precsim(y-b)_+$, then for every $\varepsilon>0$ there is $r\in D$
such that
\[
 \|u-r^*(y-a)_+r\|<\varepsilon,
 \qquad
 \|r\|^2\le\frac{\|u\|+\varepsilon}{b-a}.
\]
Here $(y-0)_+=y$.  In particular the bound is uniform in a family of
contractions $u$ with a common gap $b-a>0$.
\end{lemma}

\begin{proof}
Choose $s$ with $\|u-s^*(y-b)_+s\|<\varepsilon$ and put
\[
 f(t)=\left(\frac{(t-b)_+}{(t-a)_+}\right)^{1/2},
 \qquad r=f(y)s,
\]
where the ratio is zero when its denominator vanishes.  This is a
continuous function vanishing near zero, and
\[
 f(y)(y-a)_+f(y)=(y-b)_+,
 \qquad f(t)^2\le(b-a)^{-1}(t-b)_+.
\]
Consequently
\[
 r^*(y-a)_+r=s^*(y-b)_+s,
 \qquad
 \|r\|^2\le(b-a)^{-1}\|s^*(y-b)_+s\|,
\]
which proves the claim.  If the original comparison is formulated in
the stabilization, compression to the original corner gives $s\in D$.
\end{proof}

The choices $(a,b)=(c,2c)$ and $(a,b)=(0,c)$ will both be used.
Together with Lemma~\ref{prelim:epsilon-test}, the estimate permits
coordinatewise comparisons with a fixed spare cut to pass to the norm
ultraproduct.  Coordinatewise Cuntz comparison without such control
does not itself supply bounded representatives.

\subsection{Domination and unitary equivalence}

\begin{theorem}[Gabe's approximate domination theorem]
\label{prelim:domination}
Let $D$ be exact and let $\Phi,\Psi:D\to E$ be nuclear completely
positive maps, with $\Psi$ of order zero.  Suppose
\[
 \Phi(a)\in\Ideal_E(\Psi(a))\qquad(a\in D_+).
\]
For every finite $\mathcal F\subset D$ and $\varepsilon>0$ there are
$y_1,\ldots,y_N\in E$ such that
\[
 \left\|\Phi(a)-\sum_{l=1}^Ny_l^*\Psi(a)y_l\right\|<\varepsilon
 \qquad(a\in\mathcal F).
\]
\end{theorem}

This is the implication from pointwise ideal containment to approximate
domination in \cite[Theorem~3.3]{Gabe}, building on the ideal-separation
method of \cite[Proposition~4.2]{KR05}.  The approximation is in norm
simultaneously on $\mathcal F$, and the codomain may be a nonexact norm
ultraproduct.  No bound on the individual $y_l$ is asserted.

\begin{lemma}[Relative-commutant unitary equivalence]
\label{prelim:unitary}
Let $D_n$ be separable unital $\Z$-stable algebras, put
$D_\omega=\prod_\omega D_n$, and let $A$ be separable and unital.
For a c.p.c. order-zero map $\rho:A\to D_\omega$ set
\[
 C=D_\omega\cap\rho(A)'\cap\{1-\rho(1)\}^{\perp}.
\]
Assume that $C$ is full in $D_\omega$ and that every full hereditary
subalgebra $E\subset C$ has the following property:
whenever $x\in E$ admits totally full $e_l,e_r\in E_+$ with
$e_lx=xe_r=0$, there is a full $s\in E$ satisfying $sx=xs=0$.
If $h_0,h_1\in C_+$ are totally full contractions and
\[
 f(h_0)\sim f(h_1)\quad\text{in }C
 \qquad\text{for every }f\in C_0((0,1])_+,
\]
then $uh_1u^*=h_0$ for some unitary $u\in\widetilde C$.
\end{lemma}

This is \cite[Lemma~5.3]{BBSTWW}, whose proof adapts the
positive-element classification method of Robert--Santiago \cite{RS};
$\sim$ here denotes Cuntz equivalence.  We will verify the hereditary-subalgebra condition
directly.  After adjoining the ambient unit, the implementing unitary
belongs to $D_\omega$ and commutes with $\rho(A)$.  More explicitly,
if $C$ already has a unit $p_C$ different from the ambient unit, replace
a unitary $u\in C$ by $u+(1-p_C)$; otherwise use the canonical map
from the unitization to $C+\C1_{D_\omega}$.
